The two minute-unit displays differ by $7$ modulo $10$, whereas actual minute units must increase by $1$ modulo $10$. That position is faulty. Of the other three positions exactly one is faulty. If the hour-tens position were correct, the hours would cross from $1\ast$ to $0\ast$ within $31$ minutes, impossible. Hence hour tens is the second faulty position. The correct positions are hour units and minute tens. The earlier hour ends in $9$, and the later hour in $0$, so the possible hour transitions are $09$ to $10$ and $19$ to $20$. The latter has hour-tens digit $1$ initially, contrary to its being faulty. Thus the transition is $09$ to $10$. Starting with minute tens $4$, adding $31$ yields $10{:}11$ through $10{:}20$; the later minute tens is $2$, forcing $09{:}49\to10{:}20$. Both faulty positions are indeed wrong at both observations. The later digit sum is $\boxed3$.
